\documentclass[11pt]{article}

\usepackage[utf8]{inputenc}
\usepackage[margin=1in]{geometry}
\usepackage{amsmath,amssymb,mathtools,enumitem,fancyhdr,booktabs}
\usepackage[misc]{ifsym}
\usepackage{hyperref}

\setlength\textwidth{7in}
\setlength\parindent{0pt}
\oddsidemargin=-0.5cm
\textheight=665pt
\pagestyle{fancy}
\setlength{\headheight}{13.6pt}
\renewcommand{\headrulewidth}{0.1pt}
\renewcommand{\footrulewidth}{0.1pt}
\lhead{\scshape Pontificia Universidad Católica de Chile}
\rhead{\scshape IMC UC}
\cfoot{\thepage}

\newcommand{\norm}[1]{\left\lVert #1\right\rVert}
\newcommand{\abs}[1]{\left|#1\right|}
\def\vec   #1{\mbox{\boldmath $#1$}{}}
\def\ten   #1{\mbox{\boldmath $#1$}{}}
\def\mat   #1{\mbox{\boldmath $\mathsf #1$}}
\def\R{\mathbb R}
\def\dx{\mbox{\(\,\mathrm{d}x\)}}
\def\dX{\mbox{\(\,\mathrm{d}X\)}}
\def\ds{\mbox{\(\,\mathrm{d}s\)}}
\def\dA{\mbox{\(\,\mathrm{d}A\)}}
\DeclareMathOperator{\dive}{div}
\DeclareMathOperator{\Div}{Div}
\DeclareMathOperator{\tr}{tr}
\DeclareMathOperator{\Cof}{Cof}
\DeclareMathOperator{\spanop}{span}

\AtBeginDocument{\vspace*{-3.2\baselineskip}}

\begin{document}
    \begin{flushleft}
        \small
        \vspace{1pt}
        \textbf{IMT3130} \hfill
        \textbf{12/06/2026}
    \end{flushleft}
    \begin{center}
        \LARGE\textbf{Pauta Ayudantía 11}\\
        \vspace{3pt}
        \normalsize\textbf{Síntesis del curso: mecánica del continuo, inf--sup y FEM}\\
        \normalsize Ayudante: Bastián Herrera \Letter{\hspace{1pt}: \texttt{biherrera@uc.cl}}
        \rule{\textwidth}{0.05mm}
    \end{center}

\section*{Problemas y soluciones}

\begin{enumerate}[label=\textbf{\arabic*.}, leftmargin=*]

\item \textbf{Hiperelasticidad finita: cinemática, tensiones y esquema numérico.}

Sea $\Omega_0=(0,L_1)\times(0,L_2)\subset\R^2$ y
\[
    \vec\varphi(\vec X,t)=\mat A(t)\vec X,
    \qquad
    \mat A(t)=
    \begin{pmatrix}
        \lambda_1(t)&\gamma(t)\\
        \chi(t)&\lambda_2(t)
    \end{pmatrix},
    \qquad
    J(t)=\lambda_1\lambda_2-\gamma\chi>0.
\]
El material es Saint Venant--Kirchhoff, con
\[
    \Psi(\ten E)=\frac{\lambda_{\mathrm L}}2(\tr\ten E)^2+
    \mu_{\mathrm L}\ten E:\ten E,
    \qquad
    \ten E=\frac12(\ten F^\top\ten F-\ten I).
\]

\begin{enumerate}[label=(\alph*), leftmargin=*]
    \item \textbf{Enunciado.} Calcule $\ten F$, $J$, la velocidad material $\dot{\vec\varphi}$, la velocidad espacial $\vec v(\vec x,t)$ y la densidad espacial $\rho(\vec x,t)$. Verifique la identidad
    \[
        \dot J=J\tr(\dot{\ten F}\ten F^{-1}).
    \]

    \textbf{Solución.} Como $\vec\varphi(\vec X,t)=\mat A(t)\vec X$, el gradiente material es constante en espacio:
    \[
        \ten F=\nabla_{\vec X}\vec\varphi=\mat A(t),
        \qquad
        J=\det\ten F=\lambda_1\lambda_2-\gamma\chi.
    \]
    La velocidad material es la derivada temporal a $\vec X$ fijo:
    \[
        \dot{\vec\varphi}(\vec X,t)=\dot{\mat A}(t)\vec X
        =
        \begin{pmatrix}
            \dot\lambda_1X_1+\dot\gamma X_2\\
            \dot\chi X_1+\dot\lambda_2X_2
        \end{pmatrix}.
    \]
    Para escribir la velocidad como campo espacial, usamos $\vec x=\mat A\vec X$. Como $J>0$, la matriz $\mat A$ es invertible y
    \[
        \vec X=\mat A^{-1}\vec x.
    \]
    Por tanto,
    \[
        \vec v(\vec x,t)=\dot{\vec\varphi}(\vec X,t)\big|_{\vec X=\mat A^{-1}\vec x}
        =\dot{\mat A}\mat A^{-1}\vec x.
    \]
    En particular,
    \[
        \nabla_{\vec x}\vec v=\dot{\mat A}\mat A^{-1}=\dot{\ten F}\ten F^{-1}.
    \]
    La conservación de masa entre configuración material y espacial da
    \[
        \rho(\vec x,t)J(t)=\rho_0,
        \qquad
        \rho(\vec x,t)=\frac{\rho_0}{J(t)}.
    \]
    Finalmente,
    \[
        \dot J=\dot\lambda_1\lambda_2+
        \lambda_1\dot\lambda_2-
        \dot\gamma\chi-
        \gamma\dot\chi.
    \]
    Además,
    \[
        \ten F^{-1}=\frac1J
        \begin{pmatrix}
            \lambda_2&-\gamma\\
            -\chi&\lambda_1
        \end{pmatrix}.
    \]
    Luego
    \[
    \begin{aligned}
        J\tr(\dot{\ten F}\ten F^{-1})
        &=
        \tr\left[
        \begin{pmatrix}
            \dot\lambda_1&\dot\gamma\\
            \dot\chi&\dot\lambda_2
        \end{pmatrix}
        \begin{pmatrix}
            \lambda_2&-\gamma\\
            -\chi&\lambda_1
        \end{pmatrix}
        \right]\\
        &=
        \dot\lambda_1\lambda_2-
        \dot\gamma\chi-
        \dot\chi\gamma+
        \dot\lambda_2\lambda_1
        =\dot J.
    \end{aligned}
    \]

    \item \textbf{Enunciado.} Usando notación indicial, derive
    \[
        S_{IJ}=\frac{\partial\Psi}{\partial E_{IJ}},
        \qquad
        P_{iJ}=\frac{\partial\Psi}{\partial F_{iJ}}=F_{iI}S_{IJ}.
    \]
    Calcule las componentes de $\ten E$, $\ten S$ y $\ten P$ para el movimiento dado.

    \textbf{Solución.} Primero,
    \[
        E_{IJ}=\frac12(F_{kI}F_{kJ}-\delta_{IJ}).
    \]
    La energía se escribe como
    \[
        \Psi(\ten E)=\frac{\lambda_{\mathrm L}}2(E_{KK})^2+
        \mu_{\mathrm L}E_{IJ}E_{IJ}.
    \]
    Derivando respecto de $E_{IJ}$,
    \[
    \begin{aligned}
        \frac{\partial\Psi}{\partial E_{IJ}}
        &=
        \lambda_{\mathrm L}E_{KK}\frac{\partial E_{LL}}{\partial E_{IJ}}
        +\mu_{\mathrm L}\frac{\partial(E_{KL}E_{KL})}{\partial E_{IJ}}\\
        &=
        \lambda_{\mathrm L}(\tr\ten E)\delta_{IJ}
        +2\mu_{\mathrm L}E_{IJ}.
    \end{aligned}
    \]
    Por tanto,
    \[
        \ten S=\lambda_{\mathrm L}(\tr\ten E)\ten I+2\mu_{\mathrm L}\ten E.
    \]
    Para obtener $\ten P$, aplicamos regla de la cadena. Como
    \[
        E_{IJ}=\frac12(F_{kI}F_{kJ}-\delta_{IJ}),
    \]
    se tiene
    \[
        \frac{\partial E_{IJ}}{\partial F_{iK}}
        =
        \frac12(\delta_{IK}F_{iJ}+\delta_{JK}F_{iI}).
    \]
    Luego
    \[
    \begin{aligned}
        P_{iK}
        &=
        \frac{\partial\Psi}{\partial F_{iK}}
        =
        S_{IJ}\frac{\partial E_{IJ}}{\partial F_{iK}}\\
        &=
        \frac12S_{IJ}(\delta_{IK}F_{iJ}+\delta_{JK}F_{iI})\\
        &=
        \frac12S_{KJ}F_{iJ}+\frac12S_{IK}F_{iI}.
    \end{aligned}
    \]
    Como $\ten S$ es simétrica, ambos términos coinciden, y queda
    \[
        P_{iK}=F_{iI}S_{IK},
        \qquad
        \ten P=\ten F\ten S.
    \]

    Para la matriz dada,
    \[
        \ten F=
        \begin{pmatrix}
            \lambda_1&\gamma\\
            \chi&\lambda_2
        \end{pmatrix},
        \qquad
        \ten F^\top\ten F=
        \begin{pmatrix}
            \lambda_1^2+\chi^2&\lambda_1\gamma+\chi\lambda_2\\
            \lambda_1\gamma+\chi\lambda_2&\gamma^2+\lambda_2^2
        \end{pmatrix}.
    \]
    Así,
    \[
    \begin{aligned}
        E_{11}&=\frac12(\lambda_1^2+\chi^2-1),\\
        E_{22}&=\frac12(\gamma^2+\lambda_2^2-1),\\
        E_{12}=E_{21}&=\frac12(\lambda_1\gamma+\chi\lambda_2),
    \end{aligned}
    \]
    y
    \[
        \tr\ten E=\frac12(\lambda_1^2+\chi^2+\gamma^2+\lambda_2^2-2).
    \]
    Por tanto,
    \[
    \begin{aligned}
        S_{11}&=\lambda_{\mathrm L}\tr(\ten E)+2\mu_{\mathrm L}E_{11},\\
        S_{22}&=\lambda_{\mathrm L}\tr(\ten E)+2\mu_{\mathrm L}E_{22},\\
        S_{12}=S_{21}&=2\mu_{\mathrm L}E_{12}.
    \end{aligned}
    \]
    Finalmente,
    \[
        \ten P=	en F\ten S
        =
        \begin{pmatrix}
            \lambda_1S_{11}+\gamma S_{12} & \lambda_1S_{12}+\gamma S_{22}\\
            \chi S_{11}+\lambda_2S_{12} & \chi S_{12}+\lambda_2S_{22}
        \end{pmatrix}.
    \]

    \item \textbf{Enunciado.} Para el balance material
    \[
        \rho_0\ddot\varphi_i-P_{iJ,J}=\rho_0B_i,
    \]
    determine una fuerza de cuerpo $\vec B(\vec X,t)$ compatible con el movimiento afín y determine la tracción material $\vec T=\ten P\vec N$ sobre un borde de normal material $\vec N$.

    \textbf{Solución.} En este movimiento $\ten F$, $\ten E$, $\ten S$ y $\ten P$ dependen de $t$, pero no de $\vec X$. Por tanto,
    \[
        P_{iJ,J}=0.
    \]
    Además,
    \[
        \ddot{\vec\varphi}(\vec X,t)=\ddot{\mat A}(t)\vec X.
    \]
    El balance queda
    \[
        \rho_0\ddot{\vec\varphi}=\rho_0\vec B.
    \]
    Una fuerza de cuerpo compatible es entonces
    \[
        \vec B(\vec X,t)=\ddot{\mat A}(t)\vec X.
    \]
    En componentes,
    \[
        B_1=\ddot\lambda_1X_1+\ddot\gamma X_2,
        \qquad
        B_2=\ddot\chi X_1+\ddot\lambda_2X_2.
    \]
    La tracción material sobre un borde con normal material $\vec N$ es
    \[
        \vec T=\ten P\vec N.
    \]
    Por ejemplo, si $\vec N=\vec e_1$, entonces $\vec T=(P_{11},P_{21})^\top$; si $\vec N=\vec e_2$, entonces $\vec T=(P_{12},P_{22})^\top$.

    \item \textbf{Enunciado.} Plantee la formulación débil lagrangiana general para un problema hiperelástico con $\vec u=\bar{\vec u}$ en $\Gamma_D$ y $\ten P\vec N=\bar{\vec T}$ en $\Gamma_N$. Proponga un esquema FEM para resolver el problema no lineal quasiestático.

    \textbf{Solución.} En quasiestática se busca $\vec\varphi=\vec X+\vec u$ y
    \[
        \ten F(\vec u)=\ten I+\nabla_{\vec X}\vec u.
    \]
    Definimos
    \[
        V_{\bar u}=\{\vec u\in H^1(\Omega_0)^d:\vec u=\bar{\vec u}\text{ en }\Gamma_D\},
        \qquad
        V_0=\{\vec\eta\in H^1(\Omega_0)^d:\vec\eta=\vec0\text{ en }\Gamma_D\}.
    \]
    La ecuación fuerte quasiestática es
    \[
        -\Div\ten P(\ten F(\vec u))=\rho_0\vec B
        \qquad\text{en }\Omega_0.
    \]
    Multiplicando por $\vec\eta\in V_0$ e integrando por partes,
    \[
    \begin{aligned}
        \int_{\Omega_0}\ten P(\ten F(\vec u)):\nabla_{\vec X}\vec\eta\dX
        -\int_{\partial\Omega_0}(\ten P\vec N)\cdot\vec\eta\dA
        =
        \int_{\Omega_0}\rho_0\vec B\cdot\vec\eta\dX.
    \end{aligned}
    \]
    En $\Gamma_D$ se tiene $\vec\eta=0$, y en $\Gamma_N$ se usa $\ten P\vec N=\bar{\vec T}$. La formulación débil es: hallar $\vec u\in V_{\bar u}$ tal que
    \[
        \int_{\Omega_0}\ten P(\ten F(\vec u)):\nabla_{\vec X}\vec\eta\dX
        =
        \int_{\Omega_0}\rho_0\vec B\cdot\vec\eta\dX
        +
        \int_{\Gamma_N}\bar{\vec T}\cdot\vec\eta\dA
        \qquad\forall\vec\eta\in V_0.
    \]
    Para resolverlo con FEM se elige $V_h\subset H^1(\Omega_0)^d$. Si el dato Dirichlet es no homogéneo, se usa un lifting discreto o eliminación de grados prescritos. Se busca $\vec u_h\in V_{h,\bar u}$ tal que
    \[
        R(\vec u_h)[\vec\eta_h]=0
        \qquad\forall\vec\eta_h\in V_{h,0},
    \]
    donde
    \[
        R(\vec u_h)[\vec\eta_h]
        =
        \int_{\Omega_0}\ten P(\ten F(\vec u_h)):\nabla_{\vec X}\vec\eta_h\dX
        -
        \int_{\Omega_0}\rho_0\vec B\cdot\vec\eta_h\dX
        -
        \int_{\Gamma_N}\bar{\vec T}\cdot\vec\eta_h\dA.
    \]
    Como $\ten P(\ten F(\vec u_h))$ depende no linealmente de $\vec u_h$, se usa Newton: dado $\vec u_h^k$, hallar $\delta\vec u_h^k\in V_{h,0}$ tal que
    \[
        DR(\vec u_h^k)[\delta\vec u_h^k,\vec\eta_h]
        =
        -R(\vec u_h^k)[\vec\eta_h]
        \qquad\forall\vec\eta_h\in V_{h,0},
    \]
    y actualizar
    \[
        \vec u_h^{k+1}=\vec u_h^k+\delta\vec u_h^k.
    \]
\end{enumerate}

\item \textbf{Stokes--Brinkman con dato de borde no homogéneo e inf--sup.}

Sea $\Omega\subset\R^d$, $d\in\{2,3\}$, Lipschitz, acotado y conexo. Se considera
\[
\begin{aligned}
    -2\mu\dive D(\vec v)+\alpha\vec v+\nabla p&=\vec f &&\text{en }\Omega,\\
    \dive\vec v&=0 &&\text{en }\Omega,\\
    \vec v&=\vec g &&\text{en }\partial\Omega,\\
    \int_\Omega p\dx&=0,
\end{aligned}
\]
con $\mu>0$, $\alpha\ge0$ y $D(\vec v)=\frac12(\nabla\vec v+\nabla\vec v^\top)$. Se asume que existe un lifting solenoidal $\vec v_g\in H^1(\Omega)^d$ tal que $\vec v_g=\vec g$ en $\partial\Omega$ y $\dive\vec v_g=0$ en $\Omega$. Sea
\[
    V=H^1_0(\Omega)^d,
    \qquad
    Q=L^2_0(\Omega).
\]

\begin{enumerate}[label=(\alph*), leftmargin=*]
    \item \textbf{Enunciado.} Introduzca $\vec u=\vec v-\vec v_g$ y derive la formulación mixta para $(\vec u,p)\in V\times Q$.

    \textbf{Solución.} Como $\vec v_g$ tiene la misma traza que $\vec v$, la diferencia
    \[
        \vec u=\vec v-\vec v_g
    \]
    satisface $\vec u=\vec0$ en $\partial\Omega$. Por tanto $\vec u\in V$. Además,
    \[
        \dive\vec u=\dive\vec v-\dive\vec v_g=0-0=0.
    \]
    Tomamos $\vec w\in V$ y multiplicamos la ecuación de momentum. Integrando por partes los términos de difusión simétrica y presión, y usando que $\vec w=0$ en la frontera, se obtiene
    \[
        2\mu\int_\Omega D(\vec v):D(\vec w)\dx
        +\alpha\int_\Omega \vec v\cdot\vec w\dx
        -\int_\Omega p\,\dive\vec w\dx
        =
        \langle\vec f,\vec w\rangle.
    \]
    Reemplazando $\vec v=\vec u+\vec v_g$,
    \[
    \begin{aligned}
        &2\mu\int_\Omega D(\vec u):D(\vec w)\dx
        +\alpha\int_\Omega \vec u\cdot\vec w\dx
        -\int_\Omega p\,\dive\vec w\dx\\
        &\qquad=
        \langle\vec f,\vec w\rangle
        -2\mu\int_\Omega D(\vec v_g):D(\vec w)\dx
        -\alpha\int_\Omega \vec v_g\cdot\vec w\dx.
    \end{aligned}
    \]
    Definimos
    \[
    \begin{aligned}
        a(\vec z,\vec w)&=
        2\mu\int_\Omega D(\vec z):D(\vec w)\dx
        +\alpha\int_\Omega\vec z\cdot\vec w\dx,\\
        b(\vec w,q)&=-\int_\Omega q\,\dive\vec w\dx,\\
        L(\vec w)&=\langle\vec f,\vec w\rangle-a(\vec v_g,\vec w).
    \end{aligned}
    \]
    La formulación mixta queda: hallar $(\vec u,p)\in V\times Q$ tal que
    \[
    \begin{aligned}
        a(\vec u,\vec w)+b(\vec w,p)&=L(\vec w) &&\forall\vec w\in V,\\
        b(\vec u,q)&=0 &&\forall q\in Q.
    \end{aligned}
    \]

    \item \textbf{Enunciado.} Defina las formas $a$ y $b$. Pruebe continuidad de $a$ y $b$, y coercividad de $a$ usando Korn.

    \textbf{Solución.} Las formas ya fueron definidas en (a). Para $a$, por Cauchy--Schwarz,
    \[
    \begin{aligned}
        \abs{a(\vec z,\vec w)}
        &\le
        2\mu\norm{D(\vec z)}_{L^2}\norm{D(\vec w)}_{L^2}
        +\alpha\norm{\vec z}_{L^2}\norm{\vec w}_{L^2}\\
        &\le
        C\norm{\vec z}_{H^1}\norm{\vec w}_{H^1}.
    \end{aligned}
    \]
    Para $b$,
    \[
    \begin{aligned}
        \abs{b(\vec w,q)}
        &=
        \abs{\int_\Omega q\,\dive\vec w\dx}
        \le
        \norm{q}_{L^2}\norm{\dive\vec w}_{L^2}
        \le
        C\norm{q}_{L^2}\norm{\vec w}_{H^1}.
    \end{aligned}
    \]
    Como $\vec w\in H^1_0(\Omega)^d$, Korn entrega
    \[
        \norm{\vec w}_{H^1(\Omega)}
        \le
        C_K\norm{D(\vec w)}_{L^2(\Omega)}.
    \]
    Luego, usando $\mu>0$ y $\alpha\ge0$,
    \[
    \begin{aligned}
        a(\vec w,\vec w)
        &=
        2\mu\norm{D(\vec w)}_{L^2}^2
        +\alpha\norm{\vec w}_{L^2}^2\\
        &\ge
        2\mu\norm{D(\vec w)}_{L^2}^2
        \ge
        c\norm{\vec w}_{H^1}^2.
    \end{aligned}
    \]
    En particular, $a$ es coerciva en todo $V$, luego también en el núcleo solenoidal.

    \item \textbf{Enunciado.} Enuncie la condición inf--sup continua de la divergencia. Explique por qué controla la presión y pruebe unicidad para el problema homogéneo.

    \textbf{Solución.} La condición inf--sup continua es: existe $\beta>0$ tal que
    \[
        \inf_{q\in Q\setminus\{0\}}
        \sup_{\vec w\in V\setminus\{\vec0\}}
        \frac{\abs{b(\vec w,q)}}
        {\norm{\vec w}_{H^1(\Omega)}\norm{q}_{L^2(\Omega)}}
        \ge\beta.
    \]
    Equivalentemente, para todo $q\in Q$,
    \[
        \beta\norm{q}_{L^2}
        \le
        \sup_{\vec w\in V\setminus\{\vec0\}}
        \frac{\abs{b(\vec w,q)}}{\norm{\vec w}_{H^1}}.
    \]
    Esta desigualdad controla la presión porque permite estimar $\norm{p}_{L^2}$ usando la ecuación donde aparece $b(\vec w,p)$.

    Para probar unicidad, se consideran dos soluciones $(\vec u_1,p_1)$ y $(\vec u_2,p_2)$ del mismo problema, y se resta. La diferencia
    \[
        \delta\vec u=\vec u_1-\vec u_2,
        \qquad
        \delta p=p_1-p_2
    \]
    satisface el problema homogéneo
    \[
    \begin{aligned}
        a(\delta\vec u,\vec w)+b(\vec w,\delta p)&=0 &&\forall\vec w\in V,\\
        b(\delta\vec u,q)&=0 &&\forall q\in Q.
    \end{aligned}
    \]
    Tomando $\vec w=\delta\vec u$ en la primera ecuación y $q=\delta p$ en la segunda, se tiene
    \[
        b(\delta\vec u,\delta p)=0.
    \]
    Por tanto,
    \[
        a(\delta\vec u,\delta\vec u)=0.
    \]
    Por coercividad,
    \[
        \delta\vec u=\vec0.
    \]
    Reemplazando en la primera ecuación queda
    \[
        b(\vec w,\delta p)=0
        \qquad\forall\vec w\in V.
    \]
    Aplicando inf--sup,
    \[
        \beta\norm{\delta p}_{L^2}
        \le
        \sup_{\vec w\in V\setminus\{\vec0\}}
        \frac{\abs{b(\vec w,\delta p)}}{\norm{\vec w}_{H^1}}
        =0,
    \]
    luego $\delta p=0$. Así, la solución es única.

    \item \textbf{Enunciado.} Deduzca una estimación de estabilidad de la forma
    \[
        \norm{\vec u}_{H^1(\Omega)}+
\norm{p}_{L^2(\Omega)}
        \le
        C\left(\norm{\vec f}_{H^{-1}(\Omega)}+
\norm{\vec v_g}_{H^1(\Omega)}\right).
    \]

    \textbf{Solución.} Primero estimamos el funcional $L$. Por continuidad de $a$,
    \[
    \begin{aligned}
        \abs{L(\vec w)}
        &\le
        \norm{\vec f}_{H^{-1}}\norm{\vec w}_{H^1}
        +\abs{a(\vec v_g,\vec w)}\\
        &\le
        \left(\norm{\vec f}_{H^{-1}}+C\norm{\vec v_g}_{H^1}\right)
        \norm{\vec w}_{H^1}.
    \end{aligned}
    \]
    Por tanto,
    \[
        \norm{L}_{V'}
        \le
        C\left(\norm{\vec f}_{H^{-1}}+
\norm{\vec v_g}_{H^1}\right).
    \]
    Tomando $\vec w=\vec u$ en la primera ecuación y $q=p$ en la segunda,
    \[
        b(\vec u,p)=0,
    \]
    y entonces
    \[
        a(\vec u,\vec u)=L(\vec u).
    \]
    Por coercividad,
    \[
        c\norm{\vec u}_{H^1}^2
        \le
        \norm{L}_{V'}\norm{\vec u}_{H^1},
    \]
    luego
    \[
        \norm{\vec u}_{H^1}
        \le
        C\norm{L}_{V'}.
    \]
    Para la presión, de la primera ecuación,
    \[
        b(\vec w,p)=L(\vec w)-a(\vec u,\vec w).
    \]
    Usando inf--sup,
    \[
    \begin{aligned}
        \beta\norm{p}_{L^2}
        &\le
        \sup_{\vec w\in V\setminus\{\vec0\}}
        \frac{\abs{b(\vec w,p)}}{\norm{\vec w}_{H^1}}\\
        &\le
        \norm{L}_{V'}+C\norm{\vec u}_{H^1}
        \le
        C\norm{L}_{V'}.
    \end{aligned}
    \]
    Combinando las cotas,
    \[
        \norm{\vec u}_{H^1(\Omega)}+
\norm{p}_{L^2(\Omega)}
        \le
        C\left(\norm{\vec f}_{H^{-1}(\Omega)}+
\norm{\vec v_g}_{H^1(\Omega)}\right).
    \]
\end{enumerate}

\item \textbf{Advección--difusión--reacción: formulación, FEM e inf--sup.}

Sea $\varepsilon>0$ y
\[
    -\varepsilon u_\varepsilon''+u_\varepsilon'+u_\varepsilon=f,
    \qquad
    u_\varepsilon(0)=g_0,
    \qquad
    u_\varepsilon(1)=g_1.
\]
Sea
\[
    \ell(x)=g_0(1-x)+g_1x,
    \qquad
    w_\varepsilon=u_\varepsilon-\ell\in H^1_0(0,1).
\]

\begin{enumerate}[label=(\alph*), leftmargin=*]
    \item \textbf{Enunciado.} Derive la formulación débil para $w_\varepsilon$ y pruebe existencia y unicidad por Lax--Milgram para cada $\varepsilon>0$.

    \textbf{Solución.} Multiplicamos la ecuación por $v\in H^1_0(0,1)$ e integramos. El término difusivo se integra por partes:
    \[
    \begin{aligned}
        \int_0^1(-\varepsilon u_\varepsilon'')v\dx
        &=
        \varepsilon\int_0^1u_\varepsilon'v'\dx
        -\varepsilon[u_\varepsilon'v]_{0}^{1}\\
        &=
        \varepsilon\int_0^1u_\varepsilon'v'\dx,
    \end{aligned}
    \]
    porque $v(0)=v(1)=0$. Por tanto,
    \[
        \varepsilon\int_0^1u_\varepsilon'v'\dx
        +\int_0^1u_\varepsilon'v\dx
        +\int_0^1u_\varepsilon v\dx
        =
        \int_0^1fv\dx.
    \]
    Reemplazando $u_\varepsilon=w_\varepsilon+\ell$, se obtiene
    \[
        a_\varepsilon(w_\varepsilon,v)=L_\varepsilon(v)
        \qquad\forall v\in H^1_0(0,1),
    \]
    donde
    \[
    \begin{aligned}
        a_\varepsilon(w,v)&=
        \varepsilon\int_0^1w'v'\dx+
        \int_0^1w'v\dx+
        \int_0^1wv\dx,\\
        L_\varepsilon(v)&=
        \int_0^1fv\dx-a_\varepsilon(\ell,v).
    \end{aligned}
    \]
    La continuidad sigue de Cauchy--Schwarz:
    \[
        \abs{a_\varepsilon(w,v)}
        \le
        C_\varepsilon\norm{w}_{H^1}\norm{v}_{H^1}.
    \]
    Para coercividad,
    \[
    \begin{aligned}
        a_\varepsilon(w,w)
        &=
        \varepsilon\norm{w'}_{L^2}^2+
        \int_0^1w'w\dx+
        \norm{w}_{L^2}^2.
    \end{aligned}
    \]
    Como $w\in H^1_0(0,1)$,
    \[
        \int_0^1w'w\dx
        =
        \frac12[w^2]_{0}^{1}=0.
    \]
    Así,
    \[
        a_\varepsilon(w,w)=
        \varepsilon\norm{w'}_{L^2}^2+
\norm{w}_{L^2}^2
        \ge
        \min\{\varepsilon,1\}\norm{w}_{H^1}^2.
    \]
    Además $L_\varepsilon$ es continuo si $f\in L^2(0,1)$ y $\ell\in H^1(0,1)$. Por Lax--Milgram, existe un único $w_\varepsilon\in H^1_0(0,1)$.

    \item \textbf{Enunciado.} Para espacios conformes $V_h^{(r)}\subset H^1_0(0,1)$, $r=1,2$, formule Galerkin. Use Céa, interpolación y Aubin--Nitsche para justificar los órdenes esperados para $\varepsilon>0$ fijo.

    \textbf{Solución.} El método de Galerkin es: hallar $w_{\varepsilon,h}^{(r)}\in V_h^{(r)}$ tal que
    \[
        a_\varepsilon(w_{\varepsilon,h}^{(r)},v_h)=L_\varepsilon(v_h)
        \qquad\forall v_h\in V_h^{(r)}.
    \]
    Por consistencia, la solución exacta satisface la misma identidad para todo $v_h\in V_h^{(r)}$. Restando,
    \[
        a_\varepsilon(w_\varepsilon-w_{\varepsilon,h}^{(r)},v_h)=0
        \qquad\forall v_h\in V_h^{(r)}.
    \]
    Como $a_\varepsilon$ es continua y coerciva para $\varepsilon>0$ fijo, Céa da
    \[
        \norm{w_\varepsilon-w_{\varepsilon,h}^{(r)}}_{H^1}
        \le
        C_\varepsilon
        \inf_{z_h\in V_h^{(r)}}
        \norm{w_\varepsilon-z_h}_{H^1}.
    \]
    Tomando el interpolante nodal $\Pi_h^{(r)}w_\varepsilon$,
    \[
        \norm{w_\varepsilon-w_{\varepsilon,h}^{(r)}}_{H^1}
        \le
        C_\varepsilon h^r|w_\varepsilon|_{H^{r+1}}.
    \]
    Por tanto,
    \[
        \norm{w_\varepsilon-w_{\varepsilon,h}^{(r)}}_{H^1}=O(h^r)
        \qquad(\varepsilon\text{ fijo}).
    \]

    Para $L^2$, sea $e=w_\varepsilon-w_{\varepsilon,h}^{(r)}$ y considere el problema auxiliar: hallar $z\in H^1_0(0,1)$ tal que
    \[
        a_\varepsilon(\phi,z)=(\phi,e)_{L^2(0,1)}
        \qquad\forall\phi\in H^1_0(0,1).
    \]
    En forma fuerte, el problema auxiliar es
    \[
        -\varepsilon z''-z'+z=e,
        \qquad
        z(0)=z(1)=0.
    \]
    Entonces
    \[
    \begin{aligned}
        \norm{e}_{L^2}^2
        &=(e,e)_{L^2}
        =a_\varepsilon(e,z)\\
        &=a_\varepsilon(e,z-z_h)
        \qquad\forall z_h\in V_h^{(r)},
    \end{aligned}
    \]
    donde se usó la ortogonalidad de Galerkin. Usando continuidad de $a_\varepsilon$,
    \[
        \norm{e}_{L^2}^2
        \le
        C_\varepsilon\norm{e}_{H^1}\norm{z-z_h}_{H^1}.
    \]
    Si el problema auxiliar satisface la regularidad
    \[
        \norm{z}_{H^2}\le C_\varepsilon\norm{e}_{L^2},
    \]
    y tomamos un interpolante lineal de $z$, entonces
    \[
        \inf_{z_h\in V_h^{(r)}}\norm{z-z_h}_{H^1}
        \le
        Ch\norm{z}_{H^2}
        \le
        C_\varepsilon h\norm{e}_{L^2}.
    \]
    Por tanto,
    \[
        \norm{e}_{L^2}
        \le
        C_\varepsilon h\norm{e}_{H^1}.
    \]
    Combinando con la estimación en $H^1$,
    \[
        \norm{w_\varepsilon-w_{\varepsilon,h}^{(r)}}_{L^2}=O(h^{r+1}).
    \]

    \item \textbf{Enunciado.} Considere formalmente el problema de primer orden que se obtiene al poner $\varepsilon=0$. Identifique qué condición de borde es compatible y explique por qué no se pueden imponer ambos datos en general.

    \textbf{Solución.} Si $\varepsilon=0$, queda
    \[
        u_0'+u_0=f.
    \]
    Esta es una ecuación diferencial de primer orden. Como el coeficiente advectivo es positivo, el dato compatible es el dato de entrada:
    \[
        u_0(0)=g_0.
    \]
    Resolviendo por factor integrante,
    \[
        (e^xu_0)'=e^xf,
    \]
    y por tanto
    \[
        u_0(x)=e^{-x}g_0+\int_0^xe^{-(x-s)}f(s)\,ds.
    \]
    El valor en $x=1$ queda determinado por $f$ y $g_0$:
    \[
        u_0(1)=e^{-1}g_0+\int_0^1e^{-(1-s)}f(s)\,ds.
    \]
    En general este valor no coincide con $g_1$. Por eso, al perderse el término de segundo orden, ya no se pueden imponer simultáneamente ambos datos de borde salvo que satisfagan una condición de compatibilidad especial.

    \item \textbf{Enunciado.} Para el problema límite homogéneo, tome
    \[
        U=\{u\in H^1(0,1):u(0)=0\},
        \qquad
        V=L^2(0,1),
        \qquad
        a_0(u,v)=\int_0^1(u'+u)v\dx.
    \]
    Pruebe las condiciones de buena formulación asociadas al Lax--Milgram generalizado: estimación inf--sup e inyectividad. Muestre dos pares discretos mal elegidos donde falla el inf--sup discreto.

    \textbf{Solución.} Primero, $a_0$ es continua:
    \[
    \begin{aligned}
        \abs{a_0(u,v)}
        &\le
        \norm{u'+u}_{L^2}\norm{v}_{L^2}
        \le
        (\norm{u'}_{L^2}+\norm{u}_{L^2})\norm{v}_{L^2}\\
        &\le
        \sqrt2\norm{u}_{H^1}\norm{v}_{L^2}.
    \end{aligned}
    \]
    Para el inf--sup, fijamos $u\in U$, $u\ne0$, y tomamos $v=u'+u\in L^2(0,1)$. Entonces
    \[
        \sup_{v\in V\setminus\{0\}}
        \frac{a_0(u,v)}{\norm{v}_{L^2}}
        \ge
        \frac{\int_0^1(u'+u)^2\dx}{\norm{u'+u}_{L^2}}
        =
        \norm{u'+u}_{L^2}.
    \]
    Además,
    \[
    \begin{aligned}
        \norm{u'+u}_{L^2}^2
        &=
        \int_0^1(u')^2\dx+2\int_0^1u'u\dx+
        \int_0^1u^2\dx\\
        &=
        \norm{u'}_{L^2}^2+u(1)^2-u(0)^2+
\norm{u}_{L^2}^2\\
        &=
        \norm{u}_{H^1}^2+u(1)^2
        \ge
        \norm{u}_{H^1}^2,
    \end{aligned}
    \]
    porque $u(0)=0$. Por tanto,
    \[
        \inf_{u\in U\setminus\{0\}}
        \sup_{v\in V\setminus\{0\}}
        \frac{a_0(u,v)}{\norm{u}_{H^1}\norm{v}_{L^2}}
        \ge1.
    \]

    Falta la condición de inyectividad en la variable de prueba: si $\zeta\in L^2(0,1)$ satisface
    \[
        a_0(u,\zeta)=0
        \qquad\forall u\in U,
    \]
    debemos probar que $\zeta=0$. Para ello, definimos $u\in U$ como la solución de
    \[
        u'+u=\zeta,
        \qquad
        u(0)=0.
    \]
    Explícitamente,
    \[
        u(x)=\int_0^xe^{-(x-s)}\zeta(s)\,ds.
    \]
    Esta función pertenece a $U$ y satisface $u'+u=\zeta$. Entonces
    \[
        0=a_0(u,\zeta)=\int_0^1(u'+u)\zeta\dx
        =
        \int_0^1\zeta^2\dx
        =
        \norm{\zeta}_{L^2}^2.
    \]
    Por tanto $\zeta=0$.

    Un primer par discreto malo es
    \[
        U_h=\spanop\{x,x^2\},
        \qquad
        V_h=\spanop\{1\}.
    \]
    Si $u_h=8x-9x^2$, entonces
    \[
    \begin{aligned}
        a_0(u_h,1)
        &=
        \int_0^1(u_h'+u_h)\dx
        =
        u_h(1)-u_h(0)+\int_0^1u_h\dx\\
        &=
        (-1)-0+\int_0^1(8x-9x^2)\dx
        =
        -1+4-3=0.
    \end{aligned}
    \]
    Como $u_h\ne0$ y todo $v_h\in V_h$ es múltiplo de $1$, el supremo discreto para este $u_h$ es cero.

    Un segundo par malo es
    \[
        U_h=\spanop\{x\},
        \qquad
        V_h=\spanop\left\{1-\frac95x\right\}.
    \]
    Para $u_h=x$ y $v_h=1-\frac95x$,
    \[
    \begin{aligned}
        a_0(x,v_h)
        &=
        \int_0^1(1+x)\left(1-\frac95x\right)\dx\\
        &=
        \int_0^1\left(1-\frac45x-\frac95x^2\right)\dx
        =
        1-\frac25-\frac35=0.
    \end{aligned}
    \]
    Nuevamente, una función discreta no nula queda invisible para todas las funciones de prueba discretas. Por eso el inf--sup discreto falla.

    \item \textbf{Enunciado.} Discuta por qué las estimaciones para $\varepsilon>0$ fijo no entregan una garantía robusta cuando $\varepsilon$ es pequeño. Proponga una estrategia numérica razonable basada en estabilización o Petrov--Galerkin.

    \textbf{Solución.} La estimación de Céa para $\varepsilon>0$ fijo depende de las constantes de continuidad y coercividad de $a_\varepsilon$. La coercividad que probamos es
    \[
        a_\varepsilon(w,w)
        \ge
        \min\{\varepsilon,1\}\norm{w}_{H^1}^2.
    \]
    Por tanto, la constante de estabilidad puede crecer como $1/\varepsilon$ cuando $\varepsilon\to0^+$. Esto significa que la teoría elíptica estándar no controla uniformemente el error cuando la difusión es muy pequeña.

    Además, cuando $\varepsilon$ es pequeño el operador dominante se parece más a
    \[
        u\mapsto u'+u
    \]
    que a un operador elíptico de segundo orden. Por eso los espacios de prueba centrados y puramente conformes pueden no capturar bien la dirección de transporte. Una estrategia razonable es usar una formulación Petrov--Galerkin, escogiendo funciones de prueba alineadas con el operador de transporte--reacción, o agregar una estabilización tipo SUPG/upwind. La idea no es cambiar la solución continua, sino estabilizar la discretización en la dirección dominante del transporte.
\end{enumerate}

\end{enumerate}

\end{document}
